\par\Needspace{6\baselineskip}\medskip
\begin{center}\hypertarget{native-225-article162}{}\pdfbookmark[2]{162.}{article-162}162.\end{center}
Im \emph{zweiten} Fall, in dem $B''$ mit $B$ zusammenfällt, fällt auch $D'$ mit diesen Punkten oder mit dem Gegenpunkt zusammen. Für $AD'-\delta$ und $A''D'-\delta''$ gilt also entweder $=0$ oder $=180^0$. Aus den Gleichungen in Artikel 143 folgt daher
\SourceLeaf{226}\[
\frac{n'r'}{n}=\pm\frac{\sin\epsilon'}{\sin\epsilon}\cdot
\frac{R\sin\delta}{\sin(z+A'D-\delta')}
\]
\[
\frac{n'r'}{n''}=\pm\frac{\sin\epsilon'}{\sin\epsilon''}\cdot
\frac{R''\sin\delta''}{\sin(z+A'D''-\delta')}
\]
\[
\begin{aligned}
&R\sin\delta\sin\epsilon''\sin(z+A'D''-\delta')\notag\\
&\qquad=PR''\sin\delta''\sin\epsilon\sin(z+A'D-\delta').
\end{aligned}
\]
Daraus ist ersichtlich, dass sich $z$ unabhängig von $Q$ allein aus $P$ bestimmen lässt (es sei denn, es gilt $A'D''=A'D$ oder $=A'D\pm180^0$, womit wir zum dritten Fall kämen). Sobald $z$ gefunden ist, ist auch $r'$ bekannt. Mithilfe der Werte der Größen $\dfrac{n'r'}{n}\mathpunct{\raisebox{.45ex}{,}}$
$\dfrac{n'r'}{n''}$ erhält man dann auch $\dfrac{n}{n'}$ und $\dfrac{n''}{n'}\mathpunct{\raisebox{.45ex}{;}}$ daraus schließlich auch $Q=2\left(\dfrac{n}{n'}+\dfrac{n''}{n'}-1\right)r'^3$. Offenbar lassen sich $P$ und $Q$ also nicht als voneinander unabhängige Angaben betrachten. Sie liefern entweder nur eine einzige Angabe oder widersprüchliche Angaben. Die Lage der Punkte $C$, $C''$ bleibt in diesem Fall beliebig, sofern man sie auf demselben Großkreis wie $C'$ wählt.

Im \emph{dritten} Fall, in dem $A'$, $B$, $B'$, $B''$ auf demselben Großkreis liegen, fallen $D$ und $D''$ jeweils mit den Punkten $B''$, $B$ oder mit deren Gegenpunkten zusammen. Durch Verbindung der Gleichungen VII, VIII, IX aus Artikel 143 folgt hieraus
\[
P=\frac{R\sin\delta\sin\epsilon''}{R''\sin\delta''\sin\epsilon}
=\frac{R\sin(l'-l)}{R''\sin(l''-l')}\mathpunct{\cdot}
\]
In diesem Fall ergibt sich der Wert von $P$ also schon aus den Angaben der Aufgabe selbst, sodass die Lage der Punkte $C$, $C'$, $C''$ unbestimmt bleibt.
